% Options for packages loaded elsewhere
% Options for packages loaded elsewhere
\PassOptionsToPackage{unicode}{hyperref}
\PassOptionsToPackage{hyphens}{url}
\PassOptionsToPackage{dvipsnames,svgnames,x11names}{xcolor}
%
\documentclass[
  letterpaper,
  DIV=11,
  numbers=noendperiod]{scrartcl}
\usepackage{xcolor}
\usepackage{amsmath,amssymb}
\setcounter{secnumdepth}{5}
\usepackage{iftex}
\ifPDFTeX
  \usepackage[T1]{fontenc}
  \usepackage[utf8]{inputenc}
  \usepackage{textcomp} % provide euro and other symbols
\else % if luatex or xetex
  \usepackage{unicode-math} % this also loads fontspec
  \defaultfontfeatures{Scale=MatchLowercase}
  \defaultfontfeatures[\rmfamily]{Ligatures=TeX,Scale=1}
\fi
\usepackage{lmodern}
\ifPDFTeX\else
  % xetex/luatex font selection
\fi
% Use upquote if available, for straight quotes in verbatim environments
\IfFileExists{upquote.sty}{\usepackage{upquote}}{}
\IfFileExists{microtype.sty}{% use microtype if available
  \usepackage[]{microtype}
  \UseMicrotypeSet[protrusion]{basicmath} % disable protrusion for tt fonts
}{}
\makeatletter
\@ifundefined{KOMAClassName}{% if non-KOMA class
  \IfFileExists{parskip.sty}{%
    \usepackage{parskip}
  }{% else
    \setlength{\parindent}{0pt}
    \setlength{\parskip}{6pt plus 2pt minus 1pt}}
}{% if KOMA class
  \KOMAoptions{parskip=half}}
\makeatother
% Make \paragraph and \subparagraph free-standing
\makeatletter
\ifx\paragraph\undefined\else
  \let\oldparagraph\paragraph
  \renewcommand{\paragraph}{
    \@ifstar
      \xxxParagraphStar
      \xxxParagraphNoStar
  }
  \newcommand{\xxxParagraphStar}[1]{\oldparagraph*{#1}\mbox{}}
  \newcommand{\xxxParagraphNoStar}[1]{\oldparagraph{#1}\mbox{}}
\fi
\ifx\subparagraph\undefined\else
  \let\oldsubparagraph\subparagraph
  \renewcommand{\subparagraph}{
    \@ifstar
      \xxxSubParagraphStar
      \xxxSubParagraphNoStar
  }
  \newcommand{\xxxSubParagraphStar}[1]{\oldsubparagraph*{#1}\mbox{}}
  \newcommand{\xxxSubParagraphNoStar}[1]{\oldsubparagraph{#1}\mbox{}}
\fi
\makeatother

\usepackage{color}
\usepackage{fancyvrb}
\newcommand{\VerbBar}{|}
\newcommand{\VERB}{\Verb[commandchars=\\\{\}]}
\DefineVerbatimEnvironment{Highlighting}{Verbatim}{commandchars=\\\{\}}
% Add ',fontsize=\small' for more characters per line
\usepackage{framed}
\definecolor{shadecolor}{RGB}{241,243,245}
\newenvironment{Shaded}{\begin{snugshade}}{\end{snugshade}}
\newcommand{\AlertTok}[1]{\textcolor[rgb]{0.68,0.00,0.00}{#1}}
\newcommand{\AnnotationTok}[1]{\textcolor[rgb]{0.37,0.37,0.37}{#1}}
\newcommand{\AttributeTok}[1]{\textcolor[rgb]{0.40,0.45,0.13}{#1}}
\newcommand{\BaseNTok}[1]{\textcolor[rgb]{0.68,0.00,0.00}{#1}}
\newcommand{\BuiltInTok}[1]{\textcolor[rgb]{0.00,0.23,0.31}{#1}}
\newcommand{\CharTok}[1]{\textcolor[rgb]{0.13,0.47,0.30}{#1}}
\newcommand{\CommentTok}[1]{\textcolor[rgb]{0.37,0.37,0.37}{#1}}
\newcommand{\CommentVarTok}[1]{\textcolor[rgb]{0.37,0.37,0.37}{\textit{#1}}}
\newcommand{\ConstantTok}[1]{\textcolor[rgb]{0.56,0.35,0.01}{#1}}
\newcommand{\ControlFlowTok}[1]{\textcolor[rgb]{0.00,0.23,0.31}{\textbf{#1}}}
\newcommand{\DataTypeTok}[1]{\textcolor[rgb]{0.68,0.00,0.00}{#1}}
\newcommand{\DecValTok}[1]{\textcolor[rgb]{0.68,0.00,0.00}{#1}}
\newcommand{\DocumentationTok}[1]{\textcolor[rgb]{0.37,0.37,0.37}{\textit{#1}}}
\newcommand{\ErrorTok}[1]{\textcolor[rgb]{0.68,0.00,0.00}{#1}}
\newcommand{\ExtensionTok}[1]{\textcolor[rgb]{0.00,0.23,0.31}{#1}}
\newcommand{\FloatTok}[1]{\textcolor[rgb]{0.68,0.00,0.00}{#1}}
\newcommand{\FunctionTok}[1]{\textcolor[rgb]{0.28,0.35,0.67}{#1}}
\newcommand{\ImportTok}[1]{\textcolor[rgb]{0.00,0.46,0.62}{#1}}
\newcommand{\InformationTok}[1]{\textcolor[rgb]{0.37,0.37,0.37}{#1}}
\newcommand{\KeywordTok}[1]{\textcolor[rgb]{0.00,0.23,0.31}{\textbf{#1}}}
\newcommand{\NormalTok}[1]{\textcolor[rgb]{0.00,0.23,0.31}{#1}}
\newcommand{\OperatorTok}[1]{\textcolor[rgb]{0.37,0.37,0.37}{#1}}
\newcommand{\OtherTok}[1]{\textcolor[rgb]{0.00,0.23,0.31}{#1}}
\newcommand{\PreprocessorTok}[1]{\textcolor[rgb]{0.68,0.00,0.00}{#1}}
\newcommand{\RegionMarkerTok}[1]{\textcolor[rgb]{0.00,0.23,0.31}{#1}}
\newcommand{\SpecialCharTok}[1]{\textcolor[rgb]{0.37,0.37,0.37}{#1}}
\newcommand{\SpecialStringTok}[1]{\textcolor[rgb]{0.13,0.47,0.30}{#1}}
\newcommand{\StringTok}[1]{\textcolor[rgb]{0.13,0.47,0.30}{#1}}
\newcommand{\VariableTok}[1]{\textcolor[rgb]{0.07,0.07,0.07}{#1}}
\newcommand{\VerbatimStringTok}[1]{\textcolor[rgb]{0.13,0.47,0.30}{#1}}
\newcommand{\WarningTok}[1]{\textcolor[rgb]{0.37,0.37,0.37}{\textit{#1}}}

\usepackage{longtable,booktabs,array}
\usepackage{calc} % for calculating minipage widths
% Correct order of tables after \paragraph or \subparagraph
\usepackage{etoolbox}
\makeatletter
\patchcmd\longtable{\par}{\if@noskipsec\mbox{}\fi\par}{}{}
\makeatother
% Allow footnotes in longtable head/foot
\IfFileExists{footnotehyper.sty}{\usepackage{footnotehyper}}{\usepackage{footnote}}
\makesavenoteenv{longtable}
\usepackage{graphicx}
\makeatletter
\newsavebox\pandoc@box
\newcommand*\pandocbounded[1]{% scales image to fit in text height/width
  \sbox\pandoc@box{#1}%
  \Gscale@div\@tempa{\textheight}{\dimexpr\ht\pandoc@box+\dp\pandoc@box\relax}%
  \Gscale@div\@tempb{\linewidth}{\wd\pandoc@box}%
  \ifdim\@tempb\p@<\@tempa\p@\let\@tempa\@tempb\fi% select the smaller of both
  \ifdim\@tempa\p@<\p@\scalebox{\@tempa}{\usebox\pandoc@box}%
  \else\usebox{\pandoc@box}%
  \fi%
}
% Set default figure placement to htbp
\def\fps@figure{htbp}
\makeatother


% definitions for citeproc citations
\NewDocumentCommand\citeproctext{}{}
\NewDocumentCommand\citeproc{mm}{%
  \begingroup\def\citeproctext{#2}\cite{#1}\endgroup}
\makeatletter
 % allow citations to break across lines
 \let\@cite@ofmt\@firstofone
 % avoid brackets around text for \cite:
 \def\@biblabel#1{}
 \def\@cite#1#2{{#1\if@tempswa , #2\fi}}
\makeatother
\newlength{\cslhangindent}
\setlength{\cslhangindent}{1.5em}
\newlength{\csllabelwidth}
\setlength{\csllabelwidth}{3em}
\newenvironment{CSLReferences}[2] % #1 hanging-indent, #2 entry-spacing
 {\begin{list}{}{%
  \setlength{\itemindent}{0pt}
  \setlength{\leftmargin}{0pt}
  \setlength{\parsep}{0pt}
  % turn on hanging indent if param 1 is 1
  \ifodd #1
   \setlength{\leftmargin}{\cslhangindent}
   \setlength{\itemindent}{-1\cslhangindent}
  \fi
  % set entry spacing
  \setlength{\itemsep}{#2\baselineskip}}}
 {\end{list}}
\usepackage{calc}
\newcommand{\CSLBlock}[1]{\hfill\break\parbox[t]{\linewidth}{\strut\ignorespaces#1\strut}}
\newcommand{\CSLLeftMargin}[1]{\parbox[t]{\csllabelwidth}{\strut#1\strut}}
\newcommand{\CSLRightInline}[1]{\parbox[t]{\linewidth - \csllabelwidth}{\strut#1\strut}}
\newcommand{\CSLIndent}[1]{\hspace{\cslhangindent}#1}



\setlength{\emergencystretch}{3em} % prevent overfull lines

\providecommand{\tightlist}{%
  \setlength{\itemsep}{0pt}\setlength{\parskip}{0pt}}



 


\KOMAoption{captions}{tableheading}
\usepackage[document]{ragged2e}
\makeatletter
\@ifpackageloaded{caption}{}{\usepackage{caption}}
\AtBeginDocument{%
\ifdefined\contentsname
  \renewcommand*\contentsname{Table of contents}
\else
  \newcommand\contentsname{Table of contents}
\fi
\ifdefined\listfigurename
  \renewcommand*\listfigurename{List of Figures}
\else
  \newcommand\listfigurename{List of Figures}
\fi
\ifdefined\listtablename
  \renewcommand*\listtablename{List of Tables}
\else
  \newcommand\listtablename{List of Tables}
\fi
\ifdefined\figurename
  \renewcommand*\figurename{Figure}
\else
  \newcommand\figurename{Figure}
\fi
\ifdefined\tablename
  \renewcommand*\tablename{Table}
\else
  \newcommand\tablename{Table}
\fi
}
\@ifpackageloaded{float}{}{\usepackage{float}}
\floatstyle{ruled}
\@ifundefined{c@chapter}{\newfloat{codelisting}{h}{lop}}{\newfloat{codelisting}{h}{lop}[chapter]}
\floatname{codelisting}{Listing}
\newcommand*\listoflistings{\listof{codelisting}{List of Listings}}
\makeatother
\makeatletter
\makeatother
\makeatletter
\@ifpackageloaded{caption}{}{\usepackage{caption}}
\@ifpackageloaded{subcaption}{}{\usepackage{subcaption}}
\makeatother
\usepackage{bookmark}
\IfFileExists{xurl.sty}{\usepackage{xurl}}{} % add URL line breaks if available
\urlstyle{same}
\hypersetup{
  pdftitle={A Bayesian Approach to Address Changes in Claims Handling},
  pdfauthor={Rajesh Sahasrabuddhe, FCAS, FCIA, MAAA},
  colorlinks=true,
  linkcolor={blue},
  filecolor={Maroon},
  citecolor={Blue},
  urlcolor={Blue},
  pdfcreator={LaTeX via pandoc}}


\title{A Bayesian Approach to Address Changes in Claims Handling}
\author{Rajesh Sahasrabuddhe, FCAS, FCIA, MAAA}
\date{2026-03-08}
\begin{document}
\maketitle
\begin{abstract}
The seminal paper by James R. Berquist and Richard E. Sherman provided
systematic methods for adjusting claim triangles to account for changes
in claims handling. Those methods highlighted the dependency of future
development on case reserve adequacy and claim settlement rates. The
Berquist-Sherman method relies on restating claim triangles, rather than
modeling the dependency. In this paper, we present a Bayesian approach
to modeling the dependency in order to considering changes in claims
handling when estimating unpaid claims.
\end{abstract}

\renewcommand*\contentsname{Table of contents}
{
\hypersetup{linkcolor=}
\setcounter{tocdepth}{3}
\tableofcontents
}

\newcommand*{\bs}{Berquist-Sherman }

\section{Introduction}\label{introduction}

\begin{quote}
It's tough to make predictions, especially about the future. 
- Yogi Berra
\end{quote}

In an actuarial analysis, we need to assess the predictive value of
historical data in predicting the quantity of interest. In Loss Reserve
Adequacy Testing: A Comprehensive, Systematic Approach (Berquist and
Sherman 1977), James R. Berquist and Richard E. Sherman described
various investigations that the actuary should undertake to understand
(potential) changes in operations that might affect the predictive value
of the data.

Berquist and Sherman propose approaches to detect such changes and then
adjust the data to address data distortions. They focus on two
particular sources of changes in the historical data:

\begin{itemize}
\tightlist
\item
  Changes in case reserve adequacy
\item
  Changes in the claims settlement rate
\end{itemize}

In each case, they propose an approach to adjust the data to consider
these changes in the estimation of unpaid claims.

The goal of the approach we present in this paper is to include the
dependency in the unpaid claim estimation model. Actuaries generally
model dependency using a linear model. This is likely to be problematic
with a claims triangle, as there are a limited number of data points
relative to the number of parameters. This could result in overfitting
or situations where the model does not properly isolate the effect of
each of the independent variables.

In the approach proposed in this paper, we fit a model for emergence
using Bayesian techniques. We will use the interpretation of the
coefficient to establish prior distributions for each coefficient and
address the modeling concerns described above.

We present our proposed model using data from Berquist-Sherman to allow
the reader to compare the proposed approach with the Berquist-Sherman
baseline. The \texttt{R} code underlying this paper is available in a
\href{https://github.com/rajesh06/resv_call_2026}{code repository}
{[}Note to reviewers, I will move this to the CAS GitHub upon
completion.{]} Conveniently, the \texttt{ChainLadder} (Gesmann et al.
2024) package includes the data used in Berquist-Sherman.

\begin{Shaded}
\begin{Highlighting}[]
\NormalTok{med\_mal }\OtherTok{\textless{}{-}}\NormalTok{ ChainLadder}\SpecialCharTok{::}\NormalTok{MedMal}
\FunctionTok{names}\NormalTok{(med\_mal) }\OtherTok{\textless{}{-}} \FunctionTok{c}\NormalTok{(}\StringTok{\textquotesingle{}rept\textquotesingle{}}\NormalTok{, }\StringTok{\textquotesingle{}paid\textquotesingle{}}\NormalTok{, }\StringTok{\textquotesingle{}case\textquotesingle{}}\NormalTok{, }\StringTok{\textquotesingle{}open\textquotesingle{}}\NormalTok{)}
\NormalTok{med\_mal}\SpecialCharTok{$}\NormalTok{avg\_case }\OtherTok{\textless{}{-}}\NormalTok{ (med\_mal}\SpecialCharTok{$}\NormalTok{rept }\SpecialCharTok{{-}}\NormalTok{ med\_mal}\SpecialCharTok{$}\NormalTok{paid) }\SpecialCharTok{/}\NormalTok{ med\_mal}\SpecialCharTok{$}\NormalTok{open}

\NormalTok{med\_mal}\SpecialCharTok{$}\NormalTok{adj\_avg\_case }\OtherTok{\textless{}{-}}\NormalTok{ ChainLadder}\SpecialCharTok{::}\FunctionTok{inflateTriangle}\NormalTok{(}\AttributeTok{Triangle =}
\NormalTok{     med\_mal}\SpecialCharTok{$}\NormalTok{avg\_case, }\AttributeTok{rate =} \FloatTok{0.15}\NormalTok{)}
\NormalTok{med\_mal}\SpecialCharTok{$}\NormalTok{adj\_rept }\OtherTok{\textless{}{-}}\NormalTok{ med\_mal}\SpecialCharTok{$}\NormalTok{adj\_avg\_case }\SpecialCharTok{*}\NormalTok{ med\_mal}\SpecialCharTok{$}\NormalTok{open }\SpecialCharTok{+}\NormalTok{ med\_mal}\SpecialCharTok{$}\NormalTok{paid}

\CommentTok{\# Add row and column names}
\NormalTok{med\_mal }\OtherTok{\textless{}{-}} \FunctionTok{sapply}\NormalTok{(}\AttributeTok{X =} \FunctionTok{names}\NormalTok{(med\_mal), }\AttributeTok{FUN =} \ControlFlowTok{function}\NormalTok{(data\_type)\{}
\NormalTok{  this\_matrix }\OtherTok{\textless{}{-}}\NormalTok{ med\_mal[[data\_type]]}
  \FunctionTok{dimnames}\NormalTok{(this\_matrix) }\OtherTok{\textless{}{-}} \FunctionTok{list}\NormalTok{(}\StringTok{\textquotesingle{}ay\textquotesingle{}} \OtherTok{=} \FunctionTok{paste0}\NormalTok{(}\StringTok{\textquotesingle{}ay\_\textquotesingle{}}\NormalTok{, }\DecValTok{1969}\SpecialCharTok{:}\DecValTok{1976}\NormalTok{),}
     \StringTok{\textquotesingle{}age\textquotesingle{}} \OtherTok{=} \FunctionTok{paste0}\NormalTok{(}\StringTok{\textquotesingle{}age\_\textquotesingle{}}\NormalTok{, }\DecValTok{12} \SpecialCharTok{*} \DecValTok{1}\SpecialCharTok{:}\DecValTok{8}\NormalTok{))}
  \FunctionTok{return}\NormalTok{(this\_matrix)}
\NormalTok{\}, }\AttributeTok{simplify =} \ConstantTok{FALSE}\NormalTok{, }\AttributeTok{USE.NAMES =} \ConstantTok{TRUE}\NormalTok{)}

\NormalTok{med\_mal}\SpecialCharTok{$}\NormalTok{adj\_rept\_df }\OtherTok{\textless{}{-}}\NormalTok{ med\_mal}\SpecialCharTok{$}\NormalTok{adj\_rept[, }\DecValTok{2}\SpecialCharTok{:}\DecValTok{8}\NormalTok{] }\SpecialCharTok{/}\NormalTok{ med\_mal}\SpecialCharTok{$}\NormalTok{adj\_rept[, }\DecValTok{1}\SpecialCharTok{:}\DecValTok{7}\NormalTok{]}

\FunctionTok{colnames}\NormalTok{(med\_mal}\SpecialCharTok{$}\NormalTok{adj\_rept\_df) }\OtherTok{\textless{}{-}} \FunctionTok{paste0}\NormalTok{(}\StringTok{\textquotesingle{}age\_\textquotesingle{}}\NormalTok{, }\DecValTok{12} \SpecialCharTok{*} \DecValTok{1}\SpecialCharTok{:}\DecValTok{7}\NormalTok{, }\StringTok{\textquotesingle{}\_\textquotesingle{}}\NormalTok{, }\DecValTok{12} \SpecialCharTok{*} \DecValTok{2}\SpecialCharTok{:}\DecValTok{8}\NormalTok{)}
 
\NormalTok{bs\_ult }\OtherTok{\textless{}{-}} \ControlFlowTok{function}\NormalTok{(rptd, paid, }\AttributeTok{row\_names =} \ConstantTok{NULL}\NormalTok{) \{}
\NormalTok{  no\_cols }\OtherTok{\textless{}{-}} \FunctionTok{dim}\NormalTok{(rptd)[}\DecValTok{2}\NormalTok{]}
\NormalTok{  df\_mat }\OtherTok{\textless{}{-}}\NormalTok{ rptd[, }\DecValTok{2}\SpecialCharTok{:}\NormalTok{no\_cols] }\SpecialCharTok{/}\NormalTok{ rptd[, }\DecValTok{1}\SpecialCharTok{:}\NormalTok{(no\_cols }\SpecialCharTok{{-}} \DecValTok{1}\NormalTok{)]}
\NormalTok{  ata }\OtherTok{\textless{}{-}} \FunctionTok{apply}\NormalTok{(}\AttributeTok{X =}\NormalTok{ df\_mat, }\AttributeTok{MARGIN =} \DecValTok{2}\NormalTok{, }\AttributeTok{FUN =}\NormalTok{ mean, }\AttributeTok{na.rm =} \ConstantTok{TRUE}\NormalTok{)}
\NormalTok{  atu }\OtherTok{\textless{}{-}} \FunctionTok{c}\NormalTok{(}\DecValTok{1}\NormalTok{, ata }\SpecialCharTok{|\textgreater{}} \FunctionTok{rev}\NormalTok{() }\SpecialCharTok{|\textgreater{}} \FunctionTok{cumprod}\NormalTok{())}
\NormalTok{  latest\_rptd }\OtherTok{\textless{}{-}} \FunctionTok{sapply}\NormalTok{(}\AttributeTok{X =} \DecValTok{1}\SpecialCharTok{:}\NormalTok{no\_cols, }\ControlFlowTok{function}\NormalTok{(x) rptd[x, no\_cols }\SpecialCharTok{+} \DecValTok{1} \SpecialCharTok{{-}}\NormalTok{ x])}
\NormalTok{  latest\_paid }\OtherTok{\textless{}{-}} \FunctionTok{sapply}\NormalTok{(}\AttributeTok{X =} \DecValTok{1}\SpecialCharTok{:}\NormalTok{no\_cols, }\ControlFlowTok{function}\NormalTok{(x) paid[x, no\_cols }\SpecialCharTok{+} \DecValTok{1} \SpecialCharTok{{-}}\NormalTok{ x])}
\NormalTok{  ultimate }\OtherTok{\textless{}{-}}\NormalTok{ latest\_rptd }\SpecialCharTok{*}\NormalTok{ atu}
\NormalTok{  matrix\_out }\OtherTok{\textless{}{-}} \FunctionTok{matrix}\NormalTok{(}\AttributeTok{data =} \FunctionTok{c}\NormalTok{(latest\_rptd, atu, ultimate, latest\_paid),}
    \AttributeTok{ncol =} \DecValTok{4}\NormalTok{, }\AttributeTok{byrow =} \ConstantTok{FALSE}\NormalTok{)}
\NormalTok{  matrix\_out }\OtherTok{\textless{}{-}} \FunctionTok{cbind}\NormalTok{(matrix\_out, matrix\_out[,}\DecValTok{1}\NormalTok{] }\SpecialCharTok{*}\NormalTok{ matrix\_out[,}\DecValTok{2}\NormalTok{] }\SpecialCharTok{{-}}
\NormalTok{      matrix\_out[,}\DecValTok{4}\NormalTok{])}
  \FunctionTok{colnames}\NormalTok{(matrix\_out) }\OtherTok{\textless{}{-}} \FunctionTok{c}\NormalTok{(}\StringTok{\textquotesingle{}rptd\textquotesingle{}}\NormalTok{, }\StringTok{\textquotesingle{}atu\textquotesingle{}}\NormalTok{, }\StringTok{\textquotesingle{}ult\textquotesingle{}}\NormalTok{, }\StringTok{\textquotesingle{}paid\textquotesingle{}}\NormalTok{, }\StringTok{\textquotesingle{}reserve\textquotesingle{}}\NormalTok{)}
  \FunctionTok{rownames}\NormalTok{(matrix\_out) }\OtherTok{\textless{}{-}}\NormalTok{ row\_names}
  \FunctionTok{return}\NormalTok{(matrix\_out)}
\NormalTok{  \}}

\FunctionTok{bs\_ult}\NormalTok{(}\AttributeTok{rptd =}\NormalTok{ med\_mal}\SpecialCharTok{$}\NormalTok{adj\_rept, }\AttributeTok{paid =}\NormalTok{ med\_mal}\SpecialCharTok{$}\NormalTok{paid,}
  \AttributeTok{row\_names =} \FunctionTok{row.names}\NormalTok{(med\_mal}\SpecialCharTok{$}\NormalTok{rept))}
\end{Highlighting}
\end{Shaded}

\begin{verbatim}
            rptd       atu       ult     paid   reserve
ay_1969 23506000 1.0000000  23506000 15815000   7691000
ay_1970 32216000 0.9787289  31530731 18983000  12547731
ay_1971 48377000 0.9430549  45622165 17707000  27915165
ay_1972 61163000 1.0045027  61438397 18518000  42920397
ay_1973 73733000 0.9288155  68484352 11292000  57192352
ay_1974 63477000 1.2448873  79021714  6267000  72754714
ay_1975 48904000 1.9096082  93387478  1565000  91822478
ay_1976 15791000 7.4582396 117773062   209000 117564062
\end{verbatim}

\section{Literature Review}\label{literature-review}

The seminal work by \bs introduced deterministic methods to adjust loss
development triangles for shifts in settlement rates and case reserve
adequacy. While their methods remain industry standards, they lack a
formal probabilistic framework for quantifying uncertainty.

\section{Methodology}\label{methodology}

We first examine the effect of changes in case reserve adequacy. \bs use
medical malpractice data to illustrate their procedure. The underlying
assumption of the \bs method is that development factors are a function
of (i) the change in paid loss, (ii) the change in the number of open
claims, and (iii) the change in case reserve adequacy, where case
reserve adequacy is measured by the ratio of average case reserves,
adjusted for inflation. \bs measure case reserve adequacy separately by
development lag. However, we believe that it's reasonable to assume that
case reserve adequacy is a calendar-year measure that affects all
accident years/lags. We can express the observed development factors as
in terms of the current average case reserves:

\hfill\break
\begin{equation}\phantomsection\label{eq-ldf1}{
d_{i, j \to j+1} = 1 + \frac{p_{i, j \to j+1}+\bar{c}_{n,j+1}\times f_{i+j} \times (1+r)^{i+1-n}\times o_{i,j+1}-\bar{c}_{n,j}\times f_{i+j-1}\times (1+r)^{i-n}\times o_{i,j}}{p_{i,j}+\bar{c}_{n,j} \times f_{i+j-1} \times (1+r)^{i-n} \times o_{i,j}}
}\end{equation}

where\\
$i \in (1, n)$ is the accident year index\\
$j \in (1, m)$ is the development lag\\
$r$ is the inflation rate\\
$\bar{c}_{i,j}$ is the average case reserve\\
$o_{i,j}$ is the number of open claims\\
$p_{i,j}$ is the number of paid loss\\
$f_a$ is the calendar year case reserve adequecy factor as the end of year $a$\\

In the \bs framework, the case reserve adequacy factors, \(f\), and the
inflation (trend) rate, \(r\), are variables, whereas the other terms
are fixed.

We propose a Bayesian hierarchical model to formalize the adjustments.
Let \(C_{i,j}\) denote the cumulative claims for accident year \(i\) and
development period \(j\). We assume:

\[C_{i,j} \sim \text{Log-Normal}(\mu_{i,j}, \sigma^2)\]

Where the parameter \(\mu_{i,j}\) is conditioned on a ``settlement
speed'' factor, allowing the model to transition smoothly between
different claims-handling regimes.

\section{Discussion}\label{discussion}

Unlike the original Berquist-Sherman approach, the Bayesian method
allows for the inclusion of prior expert judgment regarding the
magnitude of the operational change.

\section{References}\label{references}

\phantomsection\label{refs}
\begin{CSLReferences}{1}{0}
\bibitem[\citeproctext]{ref-berquist_sherman_1977}
Berquist, James R., and Richard E. Sherman. 1977. {``Loss Reserve
Adequacy Testing: A Comprehensive, Systematic Approach.''}
\emph{Proceedings of the Casualty Actuarial Society} 64: 123--59.

\bibitem[\citeproctext]{ref-ChainLadder}
Gesmann, Markus, Daniel Murphy, Yanwei (Wayne) Zhang, Alessandro
Carrato, Mario Wuthrich, Fabio Concina, and Eric Dal Moro. 2024.
\emph{ChainLadder: Statistical Methods and Models for Claims Reserving
in General Insurance}.
\url{https://CRAN.R-project.org/package=ChainLadder}.

\end{CSLReferences}




\end{document}
